\chapter{PARADISO}\label{paradiso}

\hfill\break

A differenza di Campo Alfa, su Paradiso c'era un ascensore spaziale.
Sembrava un po' diverso da quello che si librava sopra la Terra, e
qualcuno mi disse che era stato messo lì dai ruhar. Questo mi sorprese,
mi ero fatto l'idea che l'ascensore ruhar avesse subito un brutto colpo
durante la battaglia, quando i kristang avevano riconquistato il
pianeta. Forse mirare agli ascensori spaziali andava contro una delle
regole non scritte? O forse era solo una questione di praticità: le
persone che tentavano di conquistare un pianeta non volevano danneggiare
un bene prezioso e i difensori pensavano che, anche se avessero perso
quella battaglia, avrebbero potuto provare a riprendersi il pianeta un
giorno? L'ascensore ruhar percorreva un fulmine come quello sulla Terra,
ma, mentre l'ascensore thuranin aveva un cavo di energia pura, quello
progettato dai ruhar aveva un cavo fisico, più sottile di un capello
umano, al centro del fulmine. Per questo motivo, i ruhar non avevano
alcuna struttura vitale all'equatore di Paradiso: se mai quel cavo si
fosse spezzato, avrebbe potuto avvolgersi attorno al pianeta diverse
volte e precipitare con un impatto considerevole, nonostante la piccola
massa. Ci dissero che al cavo erano attaccate cariche esplosive per
tutta la sua lunghezza, per ridurre il disastro nel caso in cui si fosse
spezzato o fosse stato tagliato, ma di sicuro io non avrei voluto essere
di stanza all'equatore, avrei trascorso metà del mio tempo a guardare
con preoccupazione il cielo.

La discesa in ascensore verso Paradiso non fu diversa dal viaggio in
salita dalla Terra, solo che non c'era un cuoco dell'esercito sorridente
che serviva deliziosi cheeseburger, dovevamo di nuovo accontentarci di
pasti pronti. E quell'abitacolo aveva più finestrini, e il compartimento
dei passeggeri era tre volte buone più grande di quello dell'ascensore
sulla Terra; immagino che questo fosse progettato principalmente per uso
civile. Vedevamo bene Paradiso mentre scendevamo, e scendevamo piuttosto
in fretta, c'era molta più vibrazione di quanto ricordassi dall'altra
mia esperienza con l'ascensore spaziale. Paradiso sembrava bello, almeno
aveva un sacco di verde e blu, non i marroni morti e bruciacchiati dalle
stelle di Campo Alfa. Mentre l'ascensore si avvicinava abbastanza da
poter vedere i contorni dei villaggi, mi venne in mente che presto avrei
messo piede sul mio terzo pianeta. Il primo a essere territorio nemico,
anche se quel nemico si era già arreso ai kristang. Scendendo abbastanza
in basso da distinguere i singoli edifici, mi toccai i nuovi galloni da
sergente sulle maniche. Quel posto era la casa del mio nemico, gli
alieni che avevano attaccato la Terra. Non che avessi bisogno di
ulteriori motivazioni, ma in quel momento ero pieno di determinazione a
fare tutto il possibile per portare i ruhar fuori da quel pianeta nei
tempi previsti e mostrare ai nostri salvatori kristang che gli umani
avrebbero potuto essere incaricati di compiti più importanti in futuro.
E se qualche stronzo criceto mi avesse messo i bastoni tra le ruote,
beh, mi sarei ricordato le regole di ingaggio, ma non mi sarei fatto
cagare in testa dai criceti. Avrei ricambiato loro la cortesia che ci
avevano fatto quando si erano avvicinati di soppiatto per attaccare la
Terra e avevano distrutto gran parte della nostra capacità di
generazione elettrica.

Quando l'abitacolo dell'ascensore atterrò con un leggero colpo, ci
spinsero fuori il più in fretta possibile. Non erano trascorsi dieci
minuti da quando l'ultimo di noi aveva sgomberato la recinzione attorno
al complesso dell'ascensore, che risuonarono una serie di segnali
d'allarme e l'abitacolo risalì per andare a prendere in orbita il gruppo
successivo. Io lo guardavo, sporgendo il collo allo stremo, ancora
stupito, mentre le guardie ci portavano in fretta verso una pista
d'atterraggio dove c'era una fila dei più grandi aeroplani che avessi
mai visto. Le guardie li chiamavano ``Dumbo'', per via del loro corpo
enorme e delle ali relativamente piccole. Avevano quattro motori,
incassati all'interno delle ali vicino alle radici alari, ed erano forse
sei volte più grandi dei C-17 dell'Aeronautica degli Stati Uniti cui ero
abituato. I Dumbo necessitavano solo di una breve pista, ma non potevano
decollare in modo diretto, cosa che avrebbe richiesto troppa potenza per
risultare pratica in un dirigibile. Mentre li stavo osservando, uno si
allineò e decollò, con i motori che urlavano, e giuro che la cosa fu in
volo nel giro di meno di quattrocento metri. Lo scarico del motore era
in grado di ruotare verso il basso per agevolare il decollo, poi le ali
si occupavano di fornire portanza. Una cosa così grande, che si muoveva
a ritmo così lento nell'aria, mi fece pensare che ci dovesse essere una
stringa invisibile a sostenerla. Come le Poiane, i Dumbo facevano parte
del residuo equipaggiamento ruhar che noi umani avevamo, finché i
kristang non prendevano il controllo.

Dovevamo riunirci per plotone, ma la mia squadra era da sola. Riuscii a
tenere unito il mio gruppo di fuoco, non fu difficile, e dopo un po'
scoprii dove dovevamo andare per unirci al nostro plotone. Una volta che
i sergenti scelti ebbero fatto la conta dei componenti delle squadre,
poi dei plotoni e poi a livello della compagnia, aspettammo. E
aspettammo. E aspettammo ancora un altro po', sotto il sole caldo.
Dobbiamo aver aspettato per un'ora, mentre altri plotoni e compagnie ci
passavano davanti sui Dumbo e mentre i Dumbo atterravano, imbarcavano e
decollavano di nuovo. Vedevo Polli e Poiane girare attorno al campo
d'aviazione in lontananza, erano già una vista familiare da Campo Alfa.
La benedizione era che, di qualsiasi materiale fosse fatto l'asfalto,
era di colore marrone chiaro e non cuoceva nel calore come l'asfalto
scuro avrebbe fatto sulla Terra. Dopo un'ora, con i componenti del mio
gruppo di fuoco che si comportavano al meglio per il loro nuovo
sergente, marciai infine attraverso il campo d'aviazione, passando
accanto a un paio di Dumbo vuoti in attesa, fino a raggiungere il
nostro. Era bello vedere che l'esercito degli Stati Uniti dimostrava
un'efficienza e un'organizzazione incredibili su altri pianeti come
avveniva sulla Terra. Ci ammassammo tutti insieme su sedili scomodi
fatti di rete leggera e ci sedemmo ad aspettare di nuovo. Dalle bocche
di ventilazione nella parte anteriore dell'abitacolo usciva aria fresca,
ma quando arrivava verso di noi, nella parte posteriore, era calda e
umida come quella esterna. Sulla rampa sul retro c'era un po' di
trambusto, un capitano stava indicando qualcosa su un tablet e litigando
con un tenente. Il tenente scosse la testa di fronte al capitano e
camminò verso di noi, chiaramente contrariato. Si avvicinò abbastanza
perché potessi leggere il nome, König, sul suo cartellino, mi alzai
sull'attenti e feci il saluto: «Che succede, signore?».

«Charlie Foxtrot¹²», mormorò. Un altro tipico casino fottuto. «Dicono
che siamo sull'uccello sbagliato, la risolveremo.»

Qualunque fosse il problema, cinque minuti dopo le porte si chiusero e
iniziammo a muoverci. Tutti applaudirono. Si sparse la voce che avremmo
dovuto rilassarci e sistemarci per un volo supersonico di due ore. Il
pilota fece un annuncio quando entrammo nella velocità di crociera, che
era intorno a Mach 1,4. Pur avendo viaggiato più veloce della luce,
risalendo e discendendo wormhole e tornando in orbita su una navicella,
quella velocità m'impressionava. Poi mi colpiva che l'annuncio fosse
stato fatto da una voce umana. Mi chinai verso un secondo tenente seduto
accanto a me: «Signore, chi sta pilotando questa cosa?».

Sbadigliò: «Qualche coglione dell'Aeronautica».

«Quanto addestramento di volo hanno fatto?», chiesi lentamente, non
avevo voglia di sentire la risposta che mi aspettavo. Sulla Terra, i
piloti di C-17 facevano anni di pratica prima di qualificarsi anche solo
per il posto di copilota di destra. Non c'era speranza che quei ragazzi
avessero più di un paio di settimane di esperienza di volo. In un aereo
del tutto sconosciuto, gigantesco.

«Non c'è niente di cui preoccuparsi, sergente», disse il mio
interlocutore, anche se notai che non mi guardava in faccia. «C'è un
computer che si occupa della maggior parte del volo e un kristang in
orbita può prendere il controllo e far atterrare questa cosa da remoto,
se necessario. I ragazzi davanti sono lì per fare il caffè.»

Non risi.

\hfill\break

Il volo fu tranquillo, solo che finì in un atterraggio così brusco da
farmi pensare che il pilota fosse della Marina, non dell'Aeronautica, e
stesse cercando di far atterrare il nostro Dumbo su una portaerei. Poi
vidi il ``campo d'aviazione'', che era un campo di terra, un campo che
aveva ospitato qualche genere di coltura di recente, prima che la terra
fosse sgombrata per far posto agli aerei. Senza perdere altro tempo, ci
fecero attraversare il campo a passo di marcia fino a una fila di Poiane
in attesa. Ero eccitato, non vedevo l'ora di fare un giro sull'aereo
sconosciuto. Trenta di noi si stiparono dentro una Poiana. Avrei voluto
sedermi sulla porta, con le gambe penzoloni nell'aria, col fucile sulle
ginocchia, mentre la Poiana vagava nel paesaggio. Il copilota mi mise in
riga: «Questo uccello non è un Blackhawk, sergente, viaggiamo a circa
seicentoquaranta chilometri orari quando arriviamo in quota. Se vuole
uscire dalla porta, faccia pure, ma se la chiuda dietro».

Trovai un posto, un sergente scelto mi stava rivolgendo un ampio
sorriso, poi mi disse che stava pensando di fare la stessa cosa. Anche
il volo sulla Poiana fu tranquillo, molto più tranquillo che in
qualsiasi altro elicottero in cui mi fossi trovato, e sorvolammo
innumerevoli fattorie e boschi e laghi sparsi, finché non facemmo un
giro attorno a un villaggio e atterrammo.

\hfill\break

Quando la porta della Poiana si aprì, lo sentimmo. Avete presente il
profumo del fieno appena falciato o, se avete vissuto in periferia, di
erba appena tagliata? Quel profumo pulito, fresco, tonificante, un
profumo di natura che ti ricorda l'infanzia, le corse a piedi nudi
attraverso un campo, sotto cieli blu brillante, con grandi nuvole
bianche gonfie ammassate in alto nel cielo? Sì? Non era quello. L'odore
che emanava dal campo di\ldots{} di qualsiasi cosa mangiassero i
criceti, puzzava vagamente di\ldots{}

Baker annusò e arricciò il naso. «Cazzo. Cos'è?»

Non era proprio lì, sulla nostra faccia, era più qualcosa nel sottoscala
della nostra mente o del nostro naso, un odore che ci faceva fermare a
cercare di identificarlo. Avevo una buona memoria olfattiva, ma mi ci
volle un momento per capirlo, perché era vicino, ma non troppo.
«Parmigiano», annunciai.

Non Parmigiano della serie una grande scodella di spaghetti caldi in una
fredda notte d'inverno, con un paio di polpette succose sopra, e hai
annusato tutto il giorno quell'appetitosa salsa di pomodoro che bolliva
sul fornello e ti siedi con una fetta di pane all'aglio e un bicchiere
abbondante di vino rosso e cospargi la scodella di Parmigiano fresco e
si scioglie nella salsa -- avete presente? --, non quello. Questo era
Parmigiano della serie: hai già l'influenza o una sbronza epica e il tuo
stomaco è già sul filo del rasoio, e senti odore di Parmigiano che è
stato fuori dal frigo troppo a lungo, e sai che stai per vomitare. È il
Parmigiano di cui sentimmo l'odore, proveniente dai campi.

«Parmigiano\ldots{} cazzo, puzza come i piedi di Sanchez!», disse Chen
con disgusto.

«I miei piedi? Hai mai annusato i tuoi?»

«Io i piedi me li lavo!»

«Una volta l'anno, forse.»

«Va bene, datevi una calmata, voi due», ordinai. «Forse è una sostanza
chimica che spruzzano sui loro raccolti e sparirà. Oppure ci abitueremo
e non ce ne accorgeremo più.»

«Merda.» Chen sputò. «Sono stato con Sanchez per settimane e non mi sono
ancora abituato alla puzza dei suoi piedi. Se l'intero pianeta puzza
così, sarà un lungo dislocamento.»

\hfill\break

Il resto del nostro plotone era già lì, avevano allestito una base in
tende intorno a un fienile, la città dei criceti era a mezzo chilometro
a est. Feci sistemare il mio gruppo di fuoco e il primo sergente del
plotone mi disse di presentarmi al tenente Charles in una tenda.

Feci il saluto con eleganza: «Sergente Bishop, a rapporto come da
ordini, signore».

«Bishop, mi aspettavo la tua squadra stamattina», rispose con noncuranza
al mio saluto, sembrava distratto. Stabilire la base per un plotone su
un pianeta alieno avrebbe fatto venire il mal di testa a chiunque.

«Charlie Foxtrot durante l'imbarco sulle navicelle su Campo Alfa,
signore, i kristang ci hanno messo sull'astronave sbagliata. Siamo tutti
qui ora, signore, il sergente Agnelli ci sta sistemando.»

«Sì, l'Unef è sparsa su tutto il pianeta. Sei arrivato fin qui, questo è
ciò che conta.» Indicò una mappa spianata su un tavolo, una vera mappa,
stampata su una sorta di plastica. Riproduceva circa un decimo del
continente principale, la parte in cui ci trovavamo proprio in quel
momento. Mostrava per lo più terreni agricoli, terreni agricoli
pianeggianti, con alcune basse creste che correvano da nord a sud, fiumi
e laghi qua e là. La costa si trovava a ovest, la stazione base
dell'ascensore spaziale da qualche parte fuori dalla mappa a sud-est.
«Siamo qui in questa regione», tracciò un cerchio col dito attorno a una
zona grigia ombreggiata, «i ruhar la chiamano come cazzo la chiamano»,
indicò le denominazioni, scritte in caratteri ruhar, «ma noi la
chiamiamo Pacche-culi-stan.»

«Paccheculistan?» Risi.

Lui sorrise. «La regione a nord rispetto alla nostra è il
Pacche-spalle-stan¹³, così qui a sud le pacche devono essere sul culo,
no? Il tuo gruppo di fuoco occuperà un villaggio circa cento chilometri
a ovest di qui», indicò un puntino sulla mappa. «I criceti là sono tutti
agricoltori, duecentocinquanta se si contano le fattorie circostanti.
C'è una scuola, dei granai, un deposito per i cereali, alcune case e
nient'altro.»

Un gruppo di fuoco di quattro soldati contro più di duecento criceti? Il
calcolo delle probabilità di riuscita non mi entusiasmava. Di solito un
gruppo di fuoco non operava da solo, ma in squadra con un altro, sotto
il comando di un sergente scelto. «Signore, come si può pensare che un
solo gruppo di fuoco possa tenere sotto controllo così tanti criceti?»

«Non lo pensiamo. È un esperimento», fece una smorfia, «una brillante
idea della divisione. Abbiamo squadre di reazione rapida delle
dimensioni di una compagnia sparse in giro, con Poiane e Polli,
supportate dalla potenza di fuoco dei kristang in orbita», alzò gli
occhi al soffitto della tenda. «C'è un sacco di potenza di fuoco
disponibile se ne abbiamo bisogno, quello che la divisione vuole è
evitare di averne bisogno.»

Mi crollarono le spalle. «Cuori e menti?» Questa era l'espressione usata
dall'esercito per descrivere i tentativi di convincere la popolazione
civile a cooperare con noi, o almeno a non opporci resistenza attiva.
Avevano iniziato a usarla in Vietnam e avevano continuato, con varie
forme e terminologie, in Iraq, Afghanistan, Nigeria e in pratica in ogni
luogo in cui gli Stati Uniti d'America erano stati coinvolti dal punto
di vista militare. Avevamo costruito scuole, strade, sistemi di
distribuzione dell'acqua e dell'energia e aiutato a piantare e a
raccogliere colture. A lungo termine, era più economico e più efficace
evitare di farsi altri nemici, piuttosto che uccidere nemici.
Soprattutto perché il nemico che uccidevi aveva dei parenti, e i parenti
e, con molta probabilità, tutta la loro tribù diventavano tuoi nemici.
Se ne sapete qualcosa di storia, non vi sfuggirà che questo tipo di
campagne di cuori e menti hanno, diciamo, un bilancio in chiaroscuro. Un
mese dopo aver lasciato un'area del Terzo mondo, le scuole venivano
bruciate, il sistema idrico fatto esplodere e le linee elettriche
saccheggiate per riciclare pezzi di rame. I nostri politici dichiaravano
comunque la vittoria e passavano alla crisi successiva. Hooah.

«È più simile al soft power¹⁴, Bishop. Noi non costruiremo scuole per
questi criceti, vogliamo soltanto mantenere la pace finché non se ne
andranno dal pianeta. Pensala come la tua piccola fetta di Paradiso»,
aggiunse con un ampio sorriso.

«Signore, sono un sergente minore novellino.»

«Sei anche un contadino, sei cresciuto in una fattoria. Pure Sanchez e
Baker. Questi alieni possono anche essere criceti, ma lavorano la terra,
spero che questo vi dia una capacità di comprensione che alcuni dei
nostri ragazzi di città non avranno. Questa regione si trova al centro
della linea da evacuare fuori dai confini del pianeta, quindi c'è più
tempo perché le cose vadano male.» Mi mise una mano sulla spalla e la
strinse per rassicurarmi: «Con questi criceti non devi essere gentile,
ma neanche cercare di intimidirli. Hanno già accettato di lasciare il
pianeta, l'accordo che hanno con i kristang consente loro di continuare
a dedicarsi all'allevamento e alla raccolta fino al completamento
dell'evacuazione. I ruhar pagano i kristang per trasportare i loro
raccolti. Lo so, sembra pazzesco, ma i kristang stanno combattendo
questa guerra da molto tempo e sanno quello che fanno. I ruhar vogliono
che le spedizioni di cibo continuino a circolare, dovrebbero essere ben
motivati a collaborare con noi. Se finisci nei guai, facci un fischio
con lo zPhone e faremo entrare in campo la cavalleria aerea. Se le cose
si mettono male, i kristang possono colpire qualsiasi punto dall'orbita,
e i criceti lo sanno».

Di nuovo, non ero granché fiducioso.

«Agnelli porterà l'altra metà della squadra in un villaggio a nord
rispetto a dove sarai tu. Non siete una forza di occupazione, l'Unef
chiama queste unità squadre integrate di osservazione. La G2 vuole
sapere come va con i criceti a livello locale. Tutta la sorveglianza
satellitare e aerea nella galassia non ci dirà cosa sta davvero
accadendo a terra e una squadra di ricognizione di passaggio ogni tanto
non ci fornirà sufficienti informazioni. Abbiamo bisogno di persone a
terra, che vivano tra i criceti. C'è un set completo di guide per le
squadre integrate di osservazione sul tuo tablet. Sarebbe meglio farti
un po' di formazione, ma la divisione vuole tutti sul posto il prima
possibile, per stabilire una presenza, prima che i criceti si mettano in
testa di fare la voce grossa con l'Unef.»

«Sì, signore. Abbiamo capito.» Non doveva piacermi. È il motivo per cui
si chiamano ordini, e non suggerimenti.

\hfill\break

Quando raggiunsi l'area che il plotone stava usando come parco macchine,
la mia fiducia diminuì ancora di più. «Che diavolo è questa cosa?»,
chiesi.

«È un Crivee, sergente», annunciò il sorridente meccanico.

«Un cosa?»

«Lo chiamiamo Crivee, è come il vecchio Humvee¹⁵ sulla Terra, solo che è
realizzato con i veicoli residui dei criceti. Prendiamo qualsiasi tipo
di loro camion, ci incolliamo sopra pannelli compositi per corazza
improvvisata e il Crivee è fatto.» Sembrava orgoglioso della schifezza
dentro la quale si aspettava che saremmo andati di pattuglia. Il furgone
dei gelati di Barney sembrava più capace. Se prendeste un Suv,
aggiungeste gomme grandi, ma dall'aspetto fragile, spessi schermi sui
finestrini e pannelli multicolore fissati a casaccio qua e là, avreste
qualcosa di altrettanto brutto.

«Colla?''

«Sì. Non si può saldare, perché non è metallo. Usiamo questa colla
kristang», sollevò una pistola a spruzzo, come se ne troverebbero in
qualsiasi ferramenta sulla Terra, «applichi i pannelli sul Crivee con
questo affare, poi usi questo fantastico aggeggio kristang», sembrava un
ferro da stiro, «per fissare la colla. Roba straforte.»

«Affare? Aggeggio? Non mi convinci.»

Lui fece un'alzata di spalle. «I kristang hanno nomi per tutta questa
roba, ma sono difficili da pronunciare.»

«C'è qualcosa di più pesante?», chiesi, mentre passavo le dita con fare
scettico sul pannello della corazza artigianale, che al tatto sembrava
polistirolo. Aveva uno spessore di circa dieci centimetri, che si
assottigliava sui bordi delle porte, cosa che avrebbe potuto
rappresentare un punto debole. Se mai il polistirolo avesse un punto
forte da qualche parte.

«Abbiamo carri armati Critrak, un po' come i Bradley o gli Stryker,
perché hanno le ruote, sono veicoli addetti al trasporto del personale
militare dei criceti. Questi sono usati a livello delle compagnie. Non
si preoccupi, ho visto una demo di questa corazza», colpì il polistirolo
con le nocche, «e i proiettili normali possono soltanto ammaccarne la
superficie. Bisogna colpire nello stesso punto un paio di volte, con un
proiettile esplosivo, per penetrarla. Roba dura. Dovremmo portare un po'
di questa merda sulla Terra con noi.»

«Una demo? Ti unisci a noi in pattuglia, la prima volta che ci sparano?»

Scosse la testa e sogghignò: «Non è il mio modus operandi, sergente.
Divertitevi e non graffiate la vernice. Il soldato Ringold qui vi
mostrerà come guidarli e manterrà cariche le celle a combustibile».

\hfill\break

Dopo aver preso il nostro paio di Crivee, andammo al deposito dei
rifornimenti per ottenere tutte le attrezzature di cui avremmo avuto
bisogno per il nostro incarico di squadra integrata di osservazione. Il
``deposito di rifornimenti'' era in quello che sembrava, e puzzava, come
un ex pollaio. Per prima cosa, ricevemmo equipaggiamento personale,
uniformi extra, giubbotti antiproiettile, tutta la merda di cui ha
bisogno un soldato umano vestito di tutto punto che occupa un pianeta
alieno. Poi, munizioni e armi pesanti, che significavano missili Javelin
Bravo e la nostra vecchia arma anticarro leggera non guidata, la 84 mm
svedese At4-Cs. Eravamo là, su un pianeta alieno, dei conquistatori
cazzuti, ed eravamo equipaggiati con armi che avevo usato contro i
fanatici della milizia in Nigeria. Questo era sgradevole, a mio parere.
Nessuno chiese la mia opinione.

\hfill\break

La mattina successiva eravamo pronti a partire. Il primo sergente
Mitchell venne a salutarci. «Hai sistemato tutto, Bishop?»

«Se dicessi che è stato un November Golf», ossia un disastro nello slang
dell'esercito, «farebbe qualche differenza?»

Mitchell rise. «Il tenente non ti manderebbe fuori se non avesse totale
fiducia in te, Bishop. Così, Alfa Mike Foxtrot¹⁶ a te.» Adios figlio di
puttana, ecco cosa intendeva. Non risposi.

\hfill\break

Rombammo in direzione del villaggio che ci era stato assegnato,
lasciando sulla nostra scia una nuvola di polvere che si sollevava nel
vento e che avrebbe potuto rivelare la nostra posizione a chiunque si
fosse preoccupato di guardare. Mentre ci avvicinavamo al villaggio,
raggiungemmo una piccola altura e fermai la nostra colonna. ``Colonna''
era una parola grossa: avevamo due Crivee, io e Sanchez davanti in una
sottospecie di Suv, Baker e Chen dietro in un pick-up. Entrambi i
veicoli erano carichi di roba, tra cui quello che il plotone pensava
fosse abbastanza materiale di consumo (intendendo tutto, da cibo e
medicine alla carta igienica) da sostenere quattro uomini per un mese.
Avevamo anche un carico standard di munizioni per fucili, sia con punta
regolare che esplosiva. Inoltre, due Javelin Bravo e una coppia di At4.
Niente mortai. Se ci fossero servite armi di riserva, avremmo dovuto
chiamare la cavalleria aerea. Il primo sergente mi assicurò che avremmo
ricevuto visite regolari dal quartier generale del plotone e
rifornimenti ogni due settimane. Se il primo sergente era scettico
quanto me sul concetto di soft power, teneva i propri sentimenti per sé.
Per quanto mi riguarda, davo a quell'esperimento meno di un mese prima
che la divisione, o il quartier generale dell'Unef, tornasse in sé.
Avere una parte delle nostre forze sparutamente sparpagliate in tutto il
pianeta invitava i criceti a sconfiggerci in dettaglio, il che significa
che ogni piccola unità avrebbe potuto essere sopraffatta a poco a poco.
Se fosse accaduto, i criceti avrebbero finito per essere sbattuti fuori
dall'orbita dai kristang, ma questo non sarebbe stato di grande conforto
per gli umani morti.

Come prestabilito, una coppia di Polli in ricognizione sorvolò il
villaggio a bassa quota, ad alta velocità, per attirare l'attenzione dei
criceti residenti. I Polli avevano i pod delle armi aperti e, dopo aver
sorvolato la strada principale, ruppero la formazione e salirono, con
uno che forniva copertura a nord e uno in volo a punto fisso su di noi.
Diedi un'altra occhiata al pacchetto informativo sulle mie gambe. Foto
satellitari, rapporti di una colonna di ricognizione che era passata per
il villaggio due giorni prima e documenti del governo regionale dei
criceti. Secondo i loro registri, la popolazione attuale era di 178
criceti, rispetto alla stima di 250 che avevo avuto dal plotone. Il
nostro arrivo non doveva essere una sorpresa, perché il governo dei
criceti riferiva di avere contattato gli abitanti del villaggio per
avvisarli del nostro avvicinamento e per sollecitare la cooperazione.
Avevamo anche cartelli plastificati, stampati in caratteri ruhar, che
dovevamo affiggere in luoghi importanti del villaggio. Si presume che i
cartelli elencassero le regole, tipo le ore del coprifuoco, i divieti di
portare armi e quello che i criceti dovevano fare per ospitarci. L'idea
di appendere manifesti mi mise a disagio, riscontravo una vaga
somiglianza con una delle cose che avrebbe fatto la mano pesante dei
nazisti o dei sovietici quando occupavano un villaggio.

Segnalai al pilota del Pollo che eravamo pronti e dissi a Sanchez di
proseguire. I nostri cazzuti Crivee entrarono nel villaggio a velocità
di crociera, coi motori elettrici quasi muti, i pneumatici
scricchiolanti su ghiaia e sterrato, rallentando non appena raggiungemmo
il primo edificio, l'inizio ufficiale dell'abitato. Non c'era molto da
vedere anche se, essendo io stesso originario di una piccola città
rurale, apprezzai il fatto che gli edifici, sia case che fienili,
fossero curati e ordinati. Le case avevano alberi da ombra, arbusti e
fiori, recinti e fienili erano ben tenuti. Mi aspettavo che i criceti si
fossero lasciati sfuggire le cose di mano, avessero trascurato la
manutenzione e altre attività, come piantare fiori, dato che stavano per
evacuare il pianeta. Mi aspettavo graffiti, tipo ``Gli umani devono
andare a casa'' o qualcosa del genere. Invece, vidi indizi del fatto che
i criceti stavano progettando di restare lì per un po', per esempio un
fienile cui stavano sostituendo il tetto. I ruhar non erano in grado di
accettare di aver perso il pianeta, o sapevano qualcosa che l'Unef non
sapeva? In ogni caso, l'avrei segnalato nel mio primo rapporto della
situazione al quartier generale del plotone.

Sanchez era del Kansas orientale e osservò che il villaggio, i campi e
la campagna gli ricordavano casa. Anche a me ricordavano la mia piccola
città natale. Le abitazioni erano curate, ma non raffinate, ognuna aveva
un fienile o un'officina sul retro. Nel luogo da dove venivo, sulle
facciate di case come quelle, ci sarebbero stati cartelli che
pubblicizzavano le centinaia di attività rurali che la gente svolgeva
come secondo lavoro per sbarcare il lunario. Lavori onesti, tipo vendita
di uova fresche, saldatura, piccole riparazioni meccaniche, taglio dei
capelli, taglio della legna, baby sitting. Mio padre pensava che si
potessero combinare i compiti, tipo taglio della legna con baby sitting.
Niente più ragazzini pigri seduti con i libri da colorare o a giocare
coi Lego, mio padre avrebbe detto: «Voi ragazzi fareste meglio a
tagliare e impilare quelle tre cataste di legna prima delle cinque, o
niente succhi di frutta per voi. E non piangete se vi colpite il piede
con l'accetta, camminate che vi passa. È solo una ferita superficiale».

Può darsi che stia esagerando un po'.

C'erano veicoli per la strada o nei campi e, sebbene fosse evidente che
erano alieni, mi ricordavano casa. Ogni veicolo aveva un parafango
ammaccato o un pannello della carrozzeria di un colore diverso. Erano
veicoli da lavoro, per gente che lavora. La gente che piace a me. Solo
che erano alieni e avevano attaccato la Terra senza provocazioni. Ora
che ci pensavo, speravo che nessuno dei due Crivee fosse stato
sequestrato da quel villaggio, sarebbe stata una situazione
imbarazzante.

Un criceto maschio stava camminando sul margine della strada davanti a
noi e fece un cenno di saluto con la mano, in quella che speravo fosse
un'attitudine amichevole, o almeno non minacciosa. Si capiva che era un
maschio, perché aveva la faccia del tutto ricoperta da una pelliccia
chiara e lanuginosa; le femmine avevano la faccia per lo più rivestita
da una peluria troppo fine per vederla. Inoltre, i maschi avevano grandi
baffi.

Indossava una tuta da lavoro. Salopette blu di un tessuto simile al
jeans, con toppe sulle ginocchia, una tenuta identica a quella usata da
mio padre per lavorare in casa. E una specie di berretto da baseball di
paglia, a tesa larga. Sul cappello c'era un logo, era quello di una
squadra sportiva di criceti o di un'azienda produttrice di semi, forse
il logo dell'azienda che aveva realizzato i trattori che avevamo visto
nei campi? Ordinai di fermarci, uscii e dissi a Sanchez di restare
fermo, con il motore acceso. Il mio fucile era sul sedile, avevo la mia
arma da fianco, ma per il resto ero disarmato. ``Cuori e menti'', dissi
a me stesso, cuori e menti. Vediamo se possiamo iniziare in modo
pacifico; avrei potuto diventare rude più tardi se necessario, ma un
inizio rude sarebbe stato difficile da superare. Inoltre, la coppia di
Polli, con i pod frementi, era in volo a punto fisso in piena vista,
pronta all'azione.

Il mio zPhone era già impostato sul ruhar, sollevai una mano e dissi
soltanto «Salve.»

Il criceto parlò nel suo apparecchio tipo zPhone. «Salve. Benvenuto a
Teskor. Teskor è il nostro villaggio.»

Così il nome del posto era Teskor, o almeno è così che suonava in
traduzione. Non eravamo stati in grado di leggere i caratteri ruhar
sulla mappa. Mi puntai un dito sul petto. «Io sono Joe Bishop.» Omisi il
mio rango, non essendo sicuro che ``sergente'' significasse qualcosa per
gli alieni.

Il criceto annuì: «Mi chiamano Lester Cornhut».

«Lester Cornhut?», chiesi sorpreso, ma giuro su Dio che questo è quello
che il traduttore mi disse all'orecchio. Ci era stato detto che i nomi
non erano tradotti, perciò i suoni di ``Teskor'' e ``Lester Cornhut'' mi
giunsero così come erano stati pronunciati, nella voce ruhar un po'
stridula.

Il criceto disse: «Sì», e sorrise.

Baker mi chiese sul canale tattico: «Ha appena detto che si chiama
Lester Cornhut, o Cornhole?».

«Cornhut», risposi, con la funzione traduttore in pausa, «c'è un suono T
alla fine.» Tutti e quattro ridemmo sotto i baffi e concordammo sul
fatto che Cornhut, ``capanna di mais'', fosse un nome eccellente per un
criceto contadino. Riaccesi il traduttore. «Lei è il capo di questo
villaggio, il capo di Teskor?»

«Niente capi, niente governo qui.» Tracciò col dito una linea ideale da
un'estremità all'altra del villaggio, credo che intendesse dire che un
posto così piccolo non aveva bisogno di nessuna forma di governo. «Sono
stato scelto per parlare con te. Collaboreremo come da istruzioni», il
mio traduttore fece il suono ronzante di quando non capiva qualcosa, «e
speriamo che la vostra permanenza a Teskor sia piacevole. Questo è un
bel posto.»

Che fosse di una specie nemica o no, Lester non c'entrava con l'attacco
alla Terra, era possibile che non avesse mai sentito parlare del nostro
pianeta. ``Cuori e menti'', mi dissi. Tra noi ci fu una breve
discussione, mi confermò di avere capito la data in cui la gente di
Teskor sarebbe partita per il viaggio verso l'ascensore spaziale e che
avremmo preso possesso di una fattoria alla periferia della città.
Lester mi disse che avevano già pulito il posto in vista del nostro
arrivo, la famiglia che aveva vissuto lì aveva perso un figlio quando i
kristang avevano riconquistato il pianeta e si era trasferita mesi prima
a vivere con i parenti. Attaccammo i manifesti, andammo dall'altra parte
della città e ritornammo, poi ci recammo al nostro comando di fortuna.
Tutto sommato, il nostro ingresso al villaggio fu deludente. Quando
arrivammo alla casa che avremmo abitato, feci un cenno di saluto ai
Polli con la mano e iniziammo a scaricare tutto dai nostri Crivee. Mi
sentivo come se fossimo bambini che giocano a fingersi una famiglia.
Inviai un messaggio a Shauna, ma non c'era connessione, forse non era
ancora atterrata.

Quella prima notte non riuscii a dormire molto, ero troppo eccitato per
il fatto di essere al comando della mia prima operazione, anche se stavo
solo guidando un gruppo di fuoco esperto in una missione di
ricognizione. Allestendo la nostra postazione di comando, mi sentivo un
vero sergente, un capo. Presi il primo turno di guardia e lasciai che i
ragazzi dormissero, verso mattina mi alzai perché il sonno era
imprendibile e sollevai Baker, in modo che potesse schiacciare un altro
pisolino. Oltre a essere eccitato, ero nervoso: se i criceti fossero
stati intenzionati a compiere qualche atto ostile, era probabile che
l'avrebbero fatto durante la nostra prima notte, prevenendo il fatto che
ci sistemassimo per bene, con i sistemi di sorveglianza e i campi di
fuoco estesi attorno al comando. Per fare la guardia, avevamo sistemato
una scala per salire sul tetto della casa: ecco dove mi trovavo la
mattina presto. Il tetto aveva il vantaggio di un'eccellente visuale a
trecentosessanta gradi, aveva lo svantaggio che la persona che faceva la
guardia era del tutto esposta al fuoco dei cecchini. I nostri zPhone
avevano una caratteristica ingegnosa: se la persona cui lo zPhone era
stato assegnato moriva, se il suo cuore si fermava, il dispositivo
avvisava chiunque nelle vicinanze. Il che mi dava il grande conforto,
mentre ero esposto sul tetto, che, se un cecchino mi avesse centrato, il
mio gruppo di fuoco non sarebbe stato colto di sorpresa.

L'Unef pensava che le squadre integrate di osservazione composte da
gruppi di fuoco di quattro persone, da sole in mezzo al nulla, fossero
una grande idea. Seduto sul tetto in bella vista, nelle prime ore del
giorno, non ero così sicuro che non fossimo semplici bersagli facili. La
mattina presto, quando i pensieri di una persona possono eguagliare
l'oscurità avvolgente, mi chiesi se l'Unef sperasse che le squadre
integrate di osservazione vulnerabili sarebbero state attaccate, come
scusa per mostrare ai kristang cosa gli umani erano in grado di fare in
combattimento. Dio sa che l'esercito americano mi aveva inviato in
Nigeria in pattuglie che sembravano non avere senso, se non esporci a
ordigni esplosivi e al fuoco dei cecchini. Non era della famiglia
dell'esercito che non mi fidavo, era dei politici che davano gli ordini.
Le persone a casa sarebbero state più entusiaste di sentire parlare di
umani in combattimento, che dell'Unef che svolgeva compiti di
guarnigione, solo i primi generano titoloni. Certamente i politici
adorano i titoloni.

Eppure, era tranquillo a Teskor. Così tranquillo, così silenzioso,
sedendo lassù sul tetto. Tranquillo e buio. Con il pianeta così poco
popolato, non c'era molto inquinamento luminoso; Teskor sembrava
piuttosto buia di notte, senza semafori e con solo qualche singola luce
fioca qua e là, davanti o dietro le case. Non avevo mai vissuto
un'oscurità del genere da quando avevo lasciato il Nord del Maine
affamato di elettricità. La nostra base su Campo Alfa era illuminata da
riflettori che coprivano gran parte del cielo notturno. Lì, seduto su un
tetto nel piccolo villaggio di Teskor, sul pianeta alieno di Paradiso,
lì vidi il cielo. Vidi le stelle. La Via Lattea. Così lontano dalla
Terra, le costellazioni erano tutte spostate, ma la Via Lattea si
espandeva nel cielo in tutta la sua gloria. Era ipnotizzante, dovetti
ricordare a me stesso che non dovevo fissarla, che ero sul tetto per
fare la guardia, non per guardare le stelle. Scorsi l'orizzonte e la
strada dal villaggio con un telescopio a infrarossi, era tutto libero.
Nel cielo c'era solo un pallidissimo accenno di rosa per l'alba in
arrivo, e gli insetti cominciarono a ronzare, o cinguettare, o qualsiasi
cosa facessero gli insetti su quel pianeta. Era l'ora del mattino in
cui, sulla Terra, gli uccelli avrebbero cominciato a cantare, ma su
Paradiso non c'erano uccelli nativi. La guida diceva che la biosfera non
avrebbe consentito l'evoluzione di animali volanti diversi dagli insetti
per molti milioni di anni e, poiché adesso il pianeta era occupato da
alieni, è probabile che Paradiso non avrebbe mai avuto l'opportunità di
favorire l'evoluzione spontanea degli uccelli.

Mi sedetti sul tetto guardando il cielo a est illuminarsi, ascoltando il
lieve ronzio d'insetti invisibili e sentii nostalgia di casa. Mi mancava
la Terra. Per lo più, mi mancava la Terra per come era prima
dell'invasione aliena, prima che l'elettricità fosse un lusso, prima
che, vedendo le luci scintillanti nel cielo, pensassimo a qualcosa di
diverso dalle piogge meteoriche.

Prima del Giorno di Colombo.

\hfill\break

Mentre facevamo colazione, ricevetti messaggi da Cornpone e Shauna.
Cornpone mi diceva che la nostra vecchia squadra era stata assegnata a
una base di forza di reazione rapida a circa ottocento chilometri a nord
di Teskor. Sembrava eccitato, era un incarico strabuono. «Molto meglio»,
disse beffardo, «che badare ai criceti.» Dovetti dargli ragione. Il
gruppo di fuoco mi aveva rimpiazzato con un altro, un tipo della
Carolina del Nord, e Cornpone mi disse che quello nuovo era un grande
miglioramento rispetto a me, dal momento che aveva un vero accento del
Sud, invece della mia parlata strascicata del Downeast. Shauna riferiva
che la sua unità era di stanza in via temporanea in una base logistica a
quasi sedicimila chilometri a est. Le scrissi tre messaggi, li
cancellai, poi le inviai una semplice nota. Era probabile che non ci
saremmo visti per molto tempo, forse per sempre. Provai a rispondere in
un modo che fosse amichevole, ma non troppo, perché non si sentisse in
dovere di sforzarsi troppo di restare in contatto, sarebbe stato solo
imbarazzante.

\hfill\break

Nel terzo giorno di ``pattugliamento'' del villaggio, iniziavamo tutti a
sentirci un po' scemi. Per un giro di pattugliamento, indossavamo tutto
l'armamentario di battaglia, i nostri auricolari zPhone, telecamere e
microfoni, occhiali oscurati, e marciavamo per le strade, fucili pronti,
un dito in bilico accanto alla sicura, cercando di sembrare tosti,
mentre passavamo davanti a piccole case e fienili, campi ben curati e
alla scuola, fino ai confini del villaggio, per poi voltarci e tornare
indietro. Non era facile sembrare tosti coi criceti che se ne stavano
seduti nel portico davanti casa al mattino, sorseggiando tè e
salutandoci. O lavoravano nei campi e nei granai e si fermavano per
farci un cenno di saluto mentre passavamo a piedi. O con i figli dei
criceti che si sbracciavano eccitati, desiderosi di vedere gli strani
nuovi alieni. Noi.

Per mostrare ai criceti che il nostro gruppo di fuoco aveva rinforzi
seri, una o due volte al giorno, quella prima settimana, una coppia di
Polli e Poiane arrivava in volo a bassa quota, faceva un giro attorno al
villaggio e si dirigeva di nuovo a ovest. Il gesto intendeva anche
rinfrancare la mia unità, ma vedere gli aerei scomparire oltre
l'orizzonte occidentale non faceva che ricordarci quanto fossimo
isolati.

Il punto di svolta del nostro impegno come squadra integrata di
osservazione a Teskor ci fu la mattina del settimo giorno, dopo che
avevamo effettuato il nostro pattugliamento della notte e del mattino e
io avevo inviato un doveroso rapporto sul nulla al nostro capo plotone.
Portai Sanchez con me, per ispezionare un fienile che aveva attirato la
mia attenzione. Non è che sospettassi qualche attività insolita dei
criceti, ero solo curioso di vedere l'interno di uno dei loro fienili.
Mentre camminavamo davanti alla scuola, un gruppo di piccoli criceti
stava giocando nel cortile, dando calci a un pallone che sembrava
fabbricato con l'equivalente criceto del nastro adesivo. Sapevamo che i
ruhar non ricevevano nuovi beni di lusso spediti da oltre i confini del
pianeta. I kristang permettevano di calare dall'orbita soltanto
forniture di base, qualunque altra cosa riguardasse i criceti doveva
essere spedita in orbita, senza ritorno. È probabile che i piccoli
avessero un pallone da calcio che si era sgonfiato e avessero cercato di
aggiustarlo con del nastro adesivo.

Il fienile non era niente di speciale, era costruito con travi
metalliche e una sorta di fogli di plastica estrusa che i ruhar usavano
ovunque. L'esterno era rosso, cosa che trovai interessante; forse
dipingere granai di rosso era un'altra idea quasi universale tra le
specie aliene. In un primo momento, pensai che il fienile fosse una
specie di pollaio, finché i miei occhi non si adattarono alla luce
fioca. Allora divenne strano. Là dentro allevavano degli animali su
rastrelliere lungo i lati, e puzzava come un pollaio, o un allevamento
di maiali. La cosa strana era che ogni animale, lungo circa due metri e
largo uno, aveva un cappuccio d'argento al posto della testa, da cui si
dipartivano fili e tubi fino al muro di fronte a lui. Gli animali
riposavano in culle, notai che alcuni di quelli più lontani erano più
piccoli, forse più giovani? Quello che era davvero strano era il
silenzio, niente chiocciare di polli o grugnire di maiali, solo il
ronzio di ventilatori e pompe elettriche. E c'era il mio nuovo amico
Lester Cornhut, che veniva verso di me dall'altra parte del fienile, con
un amichevole gesto di saluto. «Salve, Joseph Bishop», gridò.

«Salve, Lester Cornhut.» Dovevo avere pronunciato bene il suo nome,
perché sorrise e mi fece un breve inchino. «Che cos'è questo posto?»,
chiesi. «Cosa sono questi», indicai quegli strani\ldots{} qualunque cosa
fossero, «animali?»

«Animali? Ah. Questi non sono animali. I ruhar non mangiano animali, da
tempo abbiamo ritenuto che fosse», Lester fece una pausa per scegliere
le parole con cura, «una barbarie.» Mi fece un sorriso rammaricato. Mi
chiesi se quello che aveva detto in ruhar non fosse perfino peggio della
parola tradotta che avevo sentito.

«A me sembrano proprio animali», disse Sanchez dietro di me. «Si
muovono.»

Il signor Cornhut si avvicinò a uno degli animali, o qualunque cosa
fossero, e gli diede un colpetto. La cosa non reagì. Eppure si muoveva a
tratti, dondolando da un lato all'altro. «Queste fonti di cibo», è
probabile che questo non fosse ben tradotto, «non hanno cervello e hanno
solo un sistema nervoso limitato. I loro corpi sono quasi del tutto di
carne, geneticamente progettati in quel modo. Un computer centrale»,
indicò il cappuccio d'argento e i fili, «fa da sistema nervoso, per
questo si muovono su schemi programmati, per stimolare la crescita
muscolare.»

«Ehm, cazzo, è forte.» Sanchez era impressionato. «Allevano solo la
carne, non l'intero animale.»

Io non sapevo se essere impressionato o terrorizzato. Quelle cose erano
senza cervello, era come allevare una bistecca invece che una mucca. Era
inquietante. Non faceva una piega, era efficiente, eppure mi turbava.
Forse era solo che non ero pronto per il futuro. Quel futuro.

Dopo che fummo tornati alla base, chiesi a Baker di scavare tra le
nostre scorte e vedere se avevamo un pallone da calcio, mi era sembrato
di averne visto uno nel ``pacchetto ricreativo'' che il plotone ci aveva
dato prima di mandarci a vivere in culo a Nettuno.

Durante il pattugliamento del mattino successivo, mi misi il pallone da
calcio sotto il braccio e, mentre passavamo di nuovo a piedi davanti
alla scuola, lo lanciai dal basso verso l'alto ai bambini criceti.
All'inizio si spaventarono e si allontanarono, ma mentre tornavamo
indietro, vidi che alcuni di loro lo stavano calciando e ci salutavano
con la mano.

Quella sera, Lester Cornhut venne in visita e chiese di me. ``Show
Bishahp'', è così che pronunciò il mio nome, è probabile che si fosse
avvicinato alla realtà più di quanto io facessi con ``Lester Cornhut''.
Intendo in base a un calcolo delle probabilità, ok? Attraverso il
traduttore, chiese se fossi lo Show Bishahp che aveva catturato un
soldato ruhar durante l'\,``azione militare'' sulla Terra, che pronunciò
``Turra''.

«Sì», annuii.

Sorrise e mise in pausa il traduttore. «Ta.» Dondolò la testa in su e in
giù. «Neh.» Scosse la testa da un lato all'altro.

«Sì.» Annuii. «No.» Scossi la testa. Poi dissi: «Ta», e dondolai la
testa in su e in giù, e «Neh», e scossi la testa da un lato all'altro.
Lui fece gli stessi gesti per ``sì'' e ``no''.

Quella fu la nostra prima vera e propria comunicazione interspecifica.
Non fosse che ci trovavamo su Paradiso soltanto perché i criceti avevano
attaccato la Terra, sarei stato emozionato.

Riavviando il traduttore, mi disse che il governatore regionale dei
criceti sarebbe stato in visita a Teskor di lì a due giorni e voleva
incontrarmi, per il tè. Davvero? Ora del tè pomeridiano con il nemico?
Avremmo mangiato sandwich al cetriolo, focaccine e crumpet, qualunque
cosa siano i crumpet? Lester era a posto per essere un ruhar: era un
contadino, ed è probabile che non fosse in alcun modo coinvolto
nell'attacco alla Terra, se mai avesse sentito parlare della Terra prima
che gli umani arrivassero su Paradiso. Ma un governatore regionale
doveva far parte del governo ruhar che aveva attaccato la Terra e ucciso
gli umani. Avevo voglia di dire a Lester che il suo governatore stronzo
avrebbe potuto dirigere la sua visita dove non splende il sole. E che io
non ero un fenomeno da baraccone che lui poteva mettere in mostra:
«Fatti avanti e guarda l'umano che ha catturato un soldato ruhar!
Promettiamo che non morde! Solo cinque dollari, compreso un souvenir
gratuito!».

``Cuori e menti'', mi dissi. Cuori e menti. Se la G2 di divisione avesse
voluto informazioni, un governatore regionale avrebbe saputo molto di
più riguardo alla situazione su Paradiso del mio amico Lester, questo
era certo. «Sarei onorato di incontrare il tuo governatore», dissi,
sperando che parlare con un sorriso finto non si traducesse in sarcasmo.

«Ti porti il tè da casa? Noi abbiamo acqua calda, ma immagino tu non
beva il nostro tè, no?», chiese Lester.

Ah, per l'amor del cielo. Si aspettavano davvero che bevessi il tè con
loro? Che cazzo, in una delle nostre confezioni di cibo doveva esserci
una bustina di tè, portarla con me non significava che sarei diventato
coccoloso e avrei cantato Kumbaya con i maledetti ruhar.